\section{From a unitary pair to a near inclusion}\label{sec:near}

We now use Proposition~\ref{prop:circle} to encode two unitaries
$U,V$ on a nonzero finite-dimensional Hilbert space $D$. The resulting
von Neumann algebras will have one-sided gap bounded by a constant
times $\|UV-VU\|$.

\subsection{Matrix substitution and the algebras}
On $L=K\otimes K\otimes D$, define
\begin{equation}\label{eq:matrix-reps}
 a(g)=\sum_k A_k(g)\otimes I_K\otimes U^k,
 \qquad
 b(g)=\sum_k I_K\otimes A_k(g)\otimes V^k
 \quad(g\in F).
\end{equation}
The series converge absolutely in norm. A spectral decomposition
$U=\sum_z zP_z$ rewrites the first expression as
$\sum_z\pi_z(g)\otimes I_K\otimes P_z$, so $a$ is a unitary
representation. Diagonalizing $V$ gives the same conclusion for $b$.
There is no requirement that their ranges commute. The weak limits
from Proposition~\ref{prop:circle} become
\begin{equation}\label{eq:matrix-limits}
 a(s_j)\xrightarrow{\WOT}t p_e\otimes I_K\otimes U,
 \qquad
 b(s_j)\xrightarrow{\WOT}t I_K\otimes p_e\otimes V.
\end{equation}
Only finite spectral sums are involved, so these limits follow
directly from \eqref{eq:circle-limit}.

Put $G=F\times F$, $S=\ell^2(G)$, and $H=S\otimes L$.
Our conventions for the commuting left and right regular
representations on $S$ are
\[
 \lambda(g)\delta_x=\delta_{gx},\qquad
 \rho(g)\delta_x=\delta_{xg^{-1}}.
\]
For $x=(x_1,x_2)\in G$, put $f(x)=a(x_1)b(x_2)$ and define the
diagonal unitary $W$ by
\[
 W(\delta_x\otimes v)=\delta_x\otimes f(x)v.
\]
Although $a$ and $b$ are representations, $f$ need not be one.
Define
\begin{equation}\label{eq:algebras}
 M=W\bigl(\lambda(G)''\otimes I_L\bigr)W^*,
 \qquad N=\bigl(\rho(G)\otimes I_L\bigr)'.
\end{equation}
Primes denote commutants. Here $\lambda(G)''\otimes I_L$ means
the algebra of all $X\otimes I_L$ with $X\in\lambda(G)''$;
it is weak operator closed, as is seen by taking matrix blocks
relative to an orthonormal basis of $L$. Thus $M,N$ are von Neumann
algebras with common identity $I_H$. Notice also that
$W^*MW\subset N$, so these examples always admit a spatial embedding.

The exactly commuting case explains the choice of $W$. If $U$ and $V$
commute, then so do the ranges of $a$ and $b$, and $f$ is a
representation of $G$. In that case
\[
 W(\lambda(g)\otimes I_L)W^*=\lambda(g)\otimes f(g)\in N
 \quad(g\in G),
\]
so $M\subset N$. When $U$ and $V$ only nearly commute, we must
control the error on the entire unit ball of $M$; estimates for the
group elements alone would not suffice.

\subsection{A norm estimate on the whole algebra}
\begin{proposition}\label{prop:near}
For the algebras in \eqref{eq:algebras},
\begin{equation}\label{eq:near-bound}
 \gamma(M,N)\leq C\|UV-VU\|,
 \qquad
 C:=\sum_{k,m,p,h\in\mathbb Z}
       c_kc_mc_pc_h|m-p|\,|h|<\infty.
\end{equation}
Moreover, for every $g\in G$,
\begin{equation}\label{eq:group-bound}
 \bigl\|W(\lambda(g)\otimes I_L)W^*
            -\lambda(g)\otimes f(g)\bigr\|
 \leq C\|UV-VU\|.
\end{equation}
The constant is independent of $D,U,V$.
\end{proposition}

\begin{proof}
Write $\nu=\|UV-VU\|$. Fix $X\in\lambda(G)''$ and denote its
scalar matrix entries by $X_{x,y}$, for $x,y\in G$. Since $X$
commutes with $\rho(G)$, these entries are invariant under common
right translation:
\begin{equation}\label{eq:right-invariance}
 X_{xq,yq}=X_{x,y}\quad(x,y,q\in G).
\end{equation}
The operator $Z=W(X\otimes I_L)W^*$ has $\B(L)$-valued blocks
\[
 Z_{x,y}=X_{x,y}a(x_1)b(x_2)b(y_2)^*a(y_1)^*.
\]
We seek an approximant $Y$ with the reordered blocks
\begin{equation}\label{eq:reordered-blocks}
 Y_{x,y}=X_{x,y}a(x_1)a(y_1)^*b(x_2)b(y_2)^*
        =X_{x,y}f(xy^{-1}).
\end{equation}
The last equality uses the representation laws for $a$ and $b$
separately. Since $(xq)(yq)^{-1}=xy^{-1}$, these blocks are
right-translation invariant by \eqref{eq:right-invariance}.
Thus, once boundedness is established,
$Y$ commutes with $\rho(G)\otimes I_L$ and belongs to $N$.

The Fourier bounds establish boundedness and control the reordering
in operator norm. On $S\otimes K\otimes K$, let $D_{k,m}$ be
the diagonal operator whose block at $x\in G$ is
$A_k(x_1)\otimes A_m(x_2)$. Then $\|D_{k,m}\|\leq c_kc_m$.
Set
\[
 B_{k,m,p,h}(X)=
 D_{k,m}(X\otimes I_K\otimes I_K)D_{h,p}^*.
\]
The $A_k$ coefficients of $a$ and $b$ act on separate copies of $K$;
only their matrix factors $U^k,V^m$ act on the same space. In the
decomposition $H=(S\otimes K\otimes K)\otimes D$, expansion of the
four factors therefore gives
\begin{align}
 Z&=\sum_{k,m,p,h} B_{k,m,p,h}(X)
                   \otimes U^kV^{m-p}U^{-h},
 \label{eq:z-expansion}\\
 Y&=\sum_{k,m,p,h} B_{k,m,p,h}(X)
                   \otimes U^{k-h}V^{m-p}.
 \label{eq:y-expansion}
\end{align}
Both sums converge absolutely because
\[
 \|B_{k,m,p,h}(X)\|
       \leq c_kc_mc_pc_h\|X\|,
 \qquad \sum_{k,m,p,h}c_kc_mc_pc_h=S_0^4<\infty.
\]
Their blocks are precisely the expressions above, which also
justifies the definition of $Y$.

For integers $a,b$, telescoping gives
\begin{equation}\label{eq:power-commutator}
 \|U^aV^b-V^bU^a\|\leq |a|\,|b|\nu.
\end{equation}
For positive powers, expand a commutator one factor at a time;
negative powers follow from
$[U^{-1},V]=-U^{-1}[U,V]U^{-1}$ and its analogue for $V^{-1}$,
where $[A,B]=AB-BA$. Consequently the two matrix factors in
\eqref{eq:z-expansion} and \eqref{eq:y-expansion} differ by at most
$|m-p|\,|h|\nu$. Hence
\begin{equation}\label{eq:reordering-estimate}
 \|Z-Y\|\leq C\nu\|X\|,
 \qquad C\leq 2S_0^2S_1^2<\infty.
\end{equation}
Since $Y\in N$ and $\|Z\|=\|X\|$, this proves
\eqref{eq:near-bound}; no bound of $1$ on $\|Y\|$ is needed.

If $X=\lambda(g)$, its nonzero entries have $x=gy$. The
representation laws for $a$ and $b$ turn
\eqref{eq:reordered-blocks} into
$Y=\lambda(g)\otimes a(g_1)b(g_2)=\lambda(g)\otimes f(g)$.
Thus \eqref{eq:group-bound} is the same estimate for this choice of $X$.
\end{proof}
